\documentclass[11pt]{article}
\usepackage[top=1.7cm, bottom=1.7cm, left=1.9cm, right=1.9cm]{geometry}
\usepackage{amsmath, amssymb}
\usepackage{booktabs}
\usepackage{array}
\usepackage{enumitem}
\usepackage{xcolor}

\setlength{\parskip}{5pt}
\setlength{\parindent}{0pt}
\renewcommand{\arraystretch}{1.15}

\setlist[enumerate]{leftmargin=*, topsep=3pt, itemsep=2pt, parsep=0pt}
\setlist[itemize]{leftmargin=*, topsep=3pt, itemsep=2pt, parsep=0pt}

\newcommand{\qhead}[2]{%
  \par\vspace{4pt}\noindent\rule{\linewidth}{0.5pt}\par\vspace{1pt}%
  \noindent\textbf{\large #1}\quad\textit{#2}\par\vspace{4pt}}

\newcommand{\shead}[1]{\par\vspace{4pt}\noindent\textbf{\normalsize #1}\par\vspace{1pt}}

\begin{document}

\begin{center}
{\Large \textbf{Custom-Axis and Unsupervised-Label Questions}} \\[2pt]
{\small for an ETF trend-following pipeline ${\sim}$ Wang Yiming's framework}
\end{center}

\shead{Universe \& data}

\textbf{7 ETFs} (US equity QQQ \& IWM; international EEM; sector XLE; credit HYG; commodity GLD; bond TLT) covering distinct asset classes. Minute price\,$+$\,volume 2009-08-01 to 2026-06-01.

\shead{Chronological train / val / test split (no walk-forward)}

\smallskip
\begin{tabular}{@{}lll@{}}
\textbf{TRAIN} & 2009-08-01 $\to$ 2021-08-01 & $\sim$12 yr; fits all model params (axis threshold, HMM, VAE, XGB) \\
\textbf{VAL}   & 2021-08-01 $\to$ 2023-08-01 & $\sim$2 yr; selects recipe (HMM config, $\delta$, etc.) \\
\textbf{TEST}  & 2023-08-01 $\to$ 2026-06-01 & $\sim$3 yr OSS; live-style simulation, reported once \\
\end{tabular}

Model is trained \emph{once} on TRAIN. At VAL/TEST the model is frozen, and the custom-axis bars + features are rebuilt online from the minute stream so the simulation matches a production deployment.

\shead{Pipeline (5 modules) and the math behind modules 1 \& 2}

\smallskip
\begin{tabular}{@{}rll@{}}
\textbf{1.} & \textbf{Custom axis} (Q1) & LogDollar dollar bars; per-asset threshold tuned on TRAIN only \\
\textbf{2.} & \textbf{Labels (unsupervised)} (Q2) & sticky $K$-state Gaussian HMM, causal filter posterior \\
3. & Features                  & 512 dims: 400 integer-diff $+$ 32 fractional-diff $+$ 40 entropy $+$ 40 freq/vol \\
4. & Dim reduction             & VAE to $K_{\mathrm{latent}}{=}16$; fit on TRAIN inner-train (purge 200, embargo 100) \\
5. & Supervised model          & XGBoost binary, early-stop on TRAIN inner-val, isotonic calibration \\
\end{tabular}

\smallskip
\textbf{LogDollar axis} (causal): for minute $t$, $x_t = \log(P_t V_t)$. Maintain a causal rolling window $w{=}1950$ minutes:
\[ z_t = \frac{x_t - \mu_t^{(w)}}{\sigma_t^{(w)}}, \qquad C_t = \sum_{j \le t} z_j \qquad\text{(running z-cumsum)} \]
Emit a bar when $\max_{j \le t} C_j - C_t \ge T$ or $C_t - \min_{j \le t} C_j \ge T$. Per-ETF threshold $T$ tuned by binary search on TRAIN minutes only to hit a target bar count $N_{\mathrm{bars}}$.

\smallskip
\textbf{HMM causal filter}: fit sticky $K$-state Gaussian HMM on TRAIN bars with observation $o_t$. The \textbf{sticky floor} is a lower bound on the diagonal of the transition matrix --- after each EM M-step we enforce
$A_{kk} \,=\, P(s_{t+1}{=}k\mid s_t{=}k) \,\ge\, \mathrm{sticky\_floor}$ (then renormalise each row).
This forces regimes to be persistent: with sticky\_floor $=0.95$, the expected dwell time in state $k$ is at least $1/(1-0.95)=20$ bars. Without it, EM happily fits very fast-switching states that are useless for trend trading.

The label uses the \emph{forward filter} (no future bars), \emph{not} the smoothed posterior:
\[ \alpha_t(k) = P(s_t{=}k,\,o_{1:t}\mid\hat\theta), \quad
   \phi_t^{\mathrm{up}} = \frac{\alpha_t(\mathrm{up})}{\sum_k \alpha_t(k)}, \quad
   y_t = \mathbb{1}\{\phi_t^{\mathrm{up}} > \tau\},\ \tau{=}0.5 \]
The same one-step filter is applied at VAL and TEST.

\smallskip
\textbf{Observation choice}. $o_t = r_t$ (\textbf{1D}, univariate $\mathcal{N}(\mu_k, \sigma_k^2)$): states differ in mean only --- EM splits by \textit{direction} ($\mu_{\mathrm{up}}{>}0$ vs $\mu_{\mathrm{down}}{<}0$). $o_t = [r_t, r_t^2]$ (\textbf{2D}): adds squared-return component, so states can also differ in \emph{magnitude} --- EM splits by \emph{volatility} ($\mu_{\mathrm{up}}[r_t^2]\gg\mu_{\mathrm{down}}[r_t^2]$, direction blurred). High-vol bars are rarer, so 2-state 2D produces strongly imbalanced labels. Detrend $o_t = r_t-\mathrm{MA}_{100}$ behaves like 1D but drift-corrected.

\smallskip
\textbf{State count $K$}. After EM, states sorted by $\mu_k[0]$; \textit{up} is highest-mean state. $K{=}3$ on $r_t$: a middle \textit{flat} state absorbs neutral bars; \{flat, down\} collapse to label $0$, pulling balance near $0.5$. $K{=}4$ over-fragments (balance $0.04$--$0.33$, worse). $y_t$ binary in all cases.

\textbf{Q1 and Q2 below} are about the two design choices that emerged as least settled --- the bar-count target for module 1 and the HMM hyperparameters for module 2.

\newpage

\qhead{Q1 \enspace Custom axis}{kurt vs peak\_frac vs $N_{\mathrm{bars}}$ --- which to optimize?}

\shead{Definitions}

Bar log-return $r_{\mathrm{bar}}$, demeaned with sample std $\sigma_r$:
\[ \mathrm{kurt} \;=\; \frac{\mathbb{E}[r_{\mathrm{bar}}^4]}{(\mathbb{E}[r_{\mathrm{bar}}^2])^2} - 3 \quad\text{(excess; Gaussian $=0$)}, \qquad
   \mathrm{peak\_frac} \;=\; P(|r_{\mathrm{bar}}| < 0.5\sigma_r) \quad\text{(Gaussian $=0.383$)} \]
$N_{\mathrm{bars}}$ is the total bars over TRAIN. Larger means more training rows; smaller means thinner bars closer to IID.

\shead{Sweep across $N_{\mathrm{bars}}$ on 7 ETFs (TRAIN 12yr, threshold tuned per cell; cell shows kurt / peak\_frac)}

\smallskip
\renewcommand{\arraystretch}{1.1}
\setlength{\tabcolsep}{5pt}
\begin{footnotesize}
\begin{tabular}{@{}lccccccc@{}}
\toprule
$N_{\mathrm{bars}}$ & \textbf{QQQ} & \textbf{IWM} & \textbf{EEM} & \textbf{XLE} & \textbf{HYG} & \textbf{TLT} & \textbf{GLD} \\
\midrule
BH (minute)  & 3587 / .67 & 1388 / .65 & 2400 / .75 & 3658 / .68 & 6875 / .82 & 1286 / .68 & 941 / .72 \\
BH (hourly)  &   21 / .60 &   33 / .58 &   32 / .66 &   65 / .61 &  165 / .70 &   19 / .59 &  19 / .66 \\
BH (daily)   &    9 / .54 &   10 / .50 &    6 / .47 &   19 / .54 &   21 / .60 &    5 / .44 &   5 / .49 \\
\midrule
LogDollar 5k &    8 / .54 &    8 / .54 &    7 / .54 &   18 / .55 &   33 / .64 &    9 / .51 &   4 / .54 \\
LogDollar 10k&    9 / .58 &   12 / .56 &   13 / .60 &   43 / .59 &   46 / .68 &    9 / .57 &   8 / .59 \\
LogDollar 20k&   21 / .61 &   21 / .60 &   28 / .65 &   60 / .62 &  129 / .72 &   18 / .60 &  16 / .64 \\
LogDollar 30k&   30 / .62 &   31 / .61 &   45 / .68 &  107 / .65 &  188 / .73 &   25 / .62 &  24 / .66 \\
LogDollar 50k&  185 / .65 &   65 / .63 &  102 / .71 &  148 / .66 &  297 / .74 &   62 / .65 &  42 / .68 \\
\midrule
RealVar 5k & 14 / .39 & 4.5 / .39 & 8.1 / .40 & 14 / .40 & 28 / .44 & 4.4 / .38 & 2.6 / .38 \\
\textbf{Volume-only 5k} & \textbf{3.2 / .46} & \textbf{5.8 / .46} & \textbf{6.0 / .47} & \textbf{11.9 / .45} & 37.7 / .67 & \textbf{3.5 / .44} & \textbf{3.0 / .46} \\
RealVar 20k & 62 / .42 & 22 / .42 & 41 / .51 & 65 / .43 & 122 / .47 & 21 / .44 & 15 / .46 \\
PriceMove 20k$^{\dagger}$ &   61 / .00 &   21 / .00 &   41 / .00 &   63 / .00 &  127 / .00 &   20 / .00 &  14 / .00 \\
Volume-only 20k    &  82 / .51 &  25 / .51 &  34 / .59 &  57 / .52 & 155 / .72 &  20 / .53 & 16 / .55 \\
\midrule
Gaussian     &    0 / .38 &    0 / .38 &    0 / .38 &    0 / .38 &    0 / .38 &    0 / .38 &   0 / .38 \\
\bottomrule
\end{tabular}
\end{footnotesize}
\setlength{\tabcolsep}{6pt}

\smallskip
\textbf{Trigger definitions}: \textbf{LogDollar}: bar when cum z-score of $\log(P_t V_t)\ge T$; \textbf{RealVar}: $\sum r^2 \ge$ thresh; \textbf{PriceMove}: $|\Delta\log P|\ge\delta_P\in[12,46]$bp; \textbf{Volume-only}: cum vol $\ge V_{\mathrm{thresh}}$. \textbf{OOS} (VAL RV@TRAIN density, kurt/peak): QQQ $.09/.39$, IWM $-.19/.35$, EEM $1.87/.41$, XLE $.51/.40$, HYG $1.50/.37$, TLT $-.06/.37$, GLD $.70/.39$ --- every kurt drops vs TRAIN.

\smallskip
\textbf{What the sweep shows}: BH at any sampling rate is wildly non-Gaussian; daily cuts $\mathrm{kurt}$ ${\sim}50{\times}$ vs minute. Inside LogDollar, both metrics worsen monotonically with $N_{\mathrm{bars}}$. \textbf{At $5\mathrm{k}$ bars}, three axis types are competitive: LogDollar (kurt $4$--$33$), RealVar (kurt $2.6$--$28$), Volume-only (kurt $3$--$38$). \textbf{Per-ETF best $\mathrm{kurt}$}: Volume wins on 5/7 ETFs; LogDollar best only for HYG. \textbf{Per-ETF best $\mathrm{peak\_frac}$}: RealVar wins everywhere ($0.38$--$0.44$, Gaussian-like). PriceMove is bimodal at $\pm\delta_P$ ($^{\dagger}\mathrm{peak\_frac}{=}0$ by construction). \textbf{Note}: at $20$k, LogDollar wins; at $5$k, Volume or RealVar wins --- the optimal axis type depends on the target $N_{\mathrm{bars}}$.

\shead{The question}

\textbf{Q1.\enspace How do you trade off $\mathrm{kurt}$, $\mathrm{peak\_frac}$, and $N_{\mathrm{bars}}$?}
\begin{itemize}
\item Do you optimise a single metric, e.g.\ minimise $\mathrm{kurt}$ subject to $N_{\mathrm{bars}} \ge N_{\min}$? If so, what is $N_{\min}$?
\item Or a composite score (e.g.\ $\mathrm{kurt} \cdot |\mathrm{peak\_frac} - 0.383|$ at fixed $N_{\mathrm{bars}}$ budget)?
\item For HYG (high-yield credit, $\mathrm{kurt}\ge 28$ at every target, every axis) and XLE (energy, $\mathrm{kurt}\ge 12$): the fat tails appear \textbf{structural} --- credit-spike / oil-shock days produce return jumps that any axis aggregates into bigger bars. Do you accept this and lean on the downstream model, switch axis type, or drop those ETFs from a Wang-style universe?
\end{itemize}

\newpage

\qhead{Q2 \enspace Unsupervised labels}{HMM-config swings label balance from 0.10 to 0.87}

\shead{Definitions}

\textbf{Notation}: $r_t$ here is the \textbf{bar log-return}, i.e.\ $\log P_{\mathrm{close},t} - \log P_{\mathrm{close},t-1}$ where $t$ indexes consecutive LogDollar bars (NOT minutes, NOT log-dollar-volume $\log(P_t V_t)$ which is only used inside the axis trigger). The HMM observation $o_t$ is constructed from $r_t$ (or transforms of it: $r_t-\mathrm{MA}_{100}$, $[r_t, r_t^2]$, etc.).

For each HMM config (states $K$, observation $o_t$, sticky floor), fit on TRAIN bars and compute causal filter posterior. \textbf{The label is always binary} (even for $K{=}3$ or $K{=}4$): we pick the state with the highest mean of the first observation dimension as the \textit{up} state, then $y_t = \mathbb{1}\{\phi_t^{\mathrm{up}} > 0.5\}$. For $K{=}3$ this collapses \{flat, down\} into class $0$. Then $\mathrm{label\_mean} = \bar y = (1/N)\sum_t y_t$ (balanced $=0.5$); $\mathrm{longest\_run} =$ max consecutive same-label bars; $\mathrm{switches} =$ count of $y_t\neq y_{t-1}$. \emph{All 20k LogDollar bars per ETF.}

\shead{HMM-config sweep on 7 ETFs (cell shows label\_mean / switches; target mean $=0.5$, more switches $=$ richer signal)}

\smallskip
\renewcommand{\arraystretch}{1.1}
\setlength{\tabcolsep}{4pt}
\begin{footnotesize}
\begin{tabular}{@{}lccccccc@{}}
\toprule
\textbf{HMM (states / obs / sticky)} & \textbf{QQQ} & \textbf{IWM} & \textbf{EEM} & \textbf{XLE} & \textbf{HYG} & \textbf{TLT} & \textbf{GLD} \\
\midrule
\textbf{3 / $[r_t,|r_t|]$ / 0.95 (best)} & \textbf{.46/8.5k} & \textbf{.48/9.0k} & \textbf{.49/8.0k} & \textbf{.48/4.9k} & \textbf{.45/8.0k} & \textbf{.50/9.8k} & \textbf{.43/8.4k} \\
2 / $[r_t,|r_t|]$ / 0.95 (Wang)    & .54/9.7k & .69/5.0k & .69/5.9k & .73/4.9k & .53/9.3k & .51/9.8k & .66/6.2k \\
3 / $r_t$ / 0.95                  & .43/2.3k & .45/2.3k & .36/2.9k & .46/2.3k & .39/2.4k & .35/2.6k & .46/3.1k \\
\midrule
2 / $r_t$ (\textbf{log $P$}) / 0.95     & .67/2.6k & .70/2.3k & .64/3.3k & .71/2.4k & .77/2.2k & .65/3.0k & .60/3.4k \\
2 / $r_t^{PV}$ (\textbf{log $PV$}) / 0.95 & .45/2.5k & .31/2.3k & .33/3.0k & .31/2.6k & .38/3.4k & .33/2.5k & .39/2.5k \\
2 / $r_t{-}\mathrm{MA}_{100}$/.95 & .67/2.6k & .70/2.4k & .65/3.3k & .71/2.4k & .23/2.2k & .66/3.0k & .61/3.4k \\
\midrule
2 / $[r_t,r_t^2]$ / 0.95          & .79/4.9k & .80/4.8k & .79/5.5k & .83/4.3k & .86/3.3k & .78/5.7k & .78/5.7k \\
2 / $[r_t,r_t^2]$ / 0.99          & .81/3.7k & .81/3.4k & .81/4.6k & .85/3.4k & .87/2.9k & .79/4.2k & .79/4.7k \\
\bottomrule
\end{tabular}
\end{footnotesize}
\setlength{\tabcolsep}{6pt}

\smallskip
\textbf{What the sweep shows}: HMM hyperparameters swing $\mathrm{label\_mean}$ from $0.07$ to $0.87$. \textbf{Winning config: 3-state $[r_t, |r_t|]$} (Wang's spec + 3 states): $\mathrm{label\_mean}\in[0.43, 0.50]$ AND $5$--$10$k switches (richest signal). 2-state $[r_t,|r_t|]$ balances bond/credit/QQQ but leaves equities at $0.66$--$0.73$. $[r_t, r_t^2]$ has $L_2$-scale $r_t^2$ that dwarfs $L_1$-scale direction so the HMM picks vol regimes; $|r_t|$ has same $L_1$ scale as $r_t$ so direction stays informative. Not shown in table (negative results): \textbf{4-state $r_t$} over-fragments to mean $0.04$--$0.33$; \textbf{2-state $\mathrm{sign}(r_t)$} balances when EM converges but degenerates on IWM/EEM/XLE; \textbf{3-state $[r_t,r_t^2]$} extreme imbalance $0.07$--$0.61$. \textbf{$K$-sweep on $[r,|r|]$} (QQQ): $K{=}2$: $0.54$, $K{=}3$: $\mathbf{0.46}$, $K{=}4$: $0.42$, $K{=}5$: $0.39$ --- monotone drop as $K$ grows; $K{=}3$ is the balance sweet spot. Per Wang's spec, imbalanced labels destroy models on retraining.

\textbf{Observation: $\log P$ vs $\log P V$} (shaded row above). Replacing $r_t = \log P_t - \log P_{t-1}$ with $r_t^{PV} = \log(P_t V_t) - \log(P_{t-1} V_{t-1})$ \emph{flips} every ETF from up-majority ($0.60$--$0.77$) to up-minority ($0.31$--$0.45$) --- it does NOT balance. Reason: $\log V$ has much larger variance than $\log P$ at the bar level, so the HMM picks up \emph{volume regimes} rather than direction; the high-$\log P V$ state captures rare high-volume bars. $\log P$ is the right observation for predicting price direction; $\log P V$ would teach the model to predict volume regimes.

\textbf{$\tau$-sweep on the winning 3-state $[r_t,|r_t|]$ config} (Wang's multi-scale spec: $\tau \in \{0.5, 0.6, 0.7, 0.8, 0.9\}$ feed 5 classifiers in an ensemble):

\smallskip
\setlength{\tabcolsep}{5pt}
\begin{footnotesize}
\begin{tabular}{@{}lccccccc@{}}
\toprule
\textbf{$\tau$} & QQQ & IWM & EEM & XLE & HYG & TLT & GLD \\
\midrule
0.50 (default) & .46 & .48 & .49 & .48 & .45 & .50 & .43 \\
0.60 & .44 & .48 & .48 & .46 & .44 & .50 & .41 \\
0.70 & .43 & .48 & .48 & .44 & .43 & .50 & .40 \\
0.80 & .42 & .47 & .47 & .40 & .43 & .49 & .39 \\
0.90 & .41 & .47 & .45 & .34 & .42 & .49 & .38 \\
\bottomrule
\end{tabular}
\end{footnotesize}
\setlength{\tabcolsep}{6pt}

\smallskip
On the winner config, $\tau\in[0.5, 0.9]$ moves $\mathrm{label\_mean}$ by $0.05$--$0.13$ per ETF --- a meaningful trend-strength dial usable in Wang's multi-classifier ensemble. Switches drop only ${\sim}10\%$ at $\tau=0.9$, so signal density is preserved.

\smallskip
\textbf{Combined Q1+Q2 recipe} --- pair the winning Q1 axis (\textbf{RealVar at $5\mathrm{k}$ bars}) with the winning Q2 HMM (\textbf{3-state $[r_t,|r_t|]$}). Result on 6 of 7 ETFs (GLD has EM-init degeneracy on RV-5k): $\mathrm{peak\_frac}\in[0.38, 0.44]$ (near-Gaussian) AND $\mathrm{label\_mean}\in[0.49, 0.53]$ --- QQQ $.53$, IWM $.51$, EEM $.50$, XLE $.49$, HYG $.51$, TLT $.50$. \textbf{OOS} (TRAIN/VAL): QQQ $.53/.50$, IWM $.51/.49$, EEM $.50/.47$, XLE $.49/.52$, HYG $.51/.48$, TLT $.50/.46$, GLD $.51/.50$ ($|\Delta|\le.04$ on all 7). Robustness: replacing RV-5k with \textbf{Vol-5k} gives similar balance ($0.46$--$0.54$ on 6/7, GLD now works at $.51$, TLT $0.38$). \emph{Cleaner bars $\Rightarrow$ cleaner HMM regimes}.


\textbf{$\tau$-shift effectiveness depends on HMM config} --- on $1$D obs ($r_t$), $\tau\in\{0.5,\dots,0.9\}$ moves $\mathrm{label\_mean}$ a LOT: QQQ K=2: $0.67\to 0.44$ ($\Delta{=}0.23$); QQQ K=3: $0.43\to 0.20$ ($\Delta{=}0.23$); GLD K=3: $0.46\to 0.07$ ($\Delta{=}0.39$). On $2$D ($[r_t,r_t^2]$): bounded $\Delta\le 0.11$ because filter posterior is concentrated near 0/1.

\smallskip
\textbf{Probing 3 more levers on the imbalanced 2-state $[r_t,r_t^2]$ config}:
\begin{itemize}
\item \textbf{$\tau$-shift} (technique 2): change label cutoff $\tau$ in $y_t = \mathbb{1}\{\phi_t^{\mathrm{up}} > \tau\}$. Sweeping $\tau\in\{0.3,\dots,0.7\}$ moves QQQ $\mathrm{label\_mean}\,0.805\to 0.778$, GLD $0.786\to 0.761$. \emph{Barely moves}: $\phi_t^{\mathrm{up}}$ is concentrated near $0$ or $1$ on 2D obs.
\item \textbf{Sticky-floor sweep}, sticky $\in\{0.50,\dots,0.99\}$: QQQ $\mathrm{label\_mean}\,0.81\to 0.73$ (switches $3.7$k $\to$ $7.1$k); GLD $0.79\to 0.75$. Moves $\sim 0.06$ at best; \emph{bounded}.
\item \textbf{Pool-shrink}, refit HMM on last $\{5,10,15,20\}$k bars only: QQQ stays $0.78$--$0.80$; GLD stays $\sim 0.78$. \emph{No systematic effect} --- imbalance is structural, not from non-stationarity.
\end{itemize}

\shead{The question}

\textbf{Q2.\enspace Should we force $\mathrm{label\_mean} \to 0.5$, and how?}

\smallskip
Levers we tested above (\emph{all unsupervised, all leak-free}):
\begin{itemize}
\item \textbf{HMM config sweep}: state count $K\in\{2,3,4\}$, observation $o_t\in\{r_t, [r_t,r_t^2], r_t-\mathrm{MA}_{100}, \mathrm{sign}(r_t)\}$, sticky $\in\{0.95, 0.99\}$. Result: only \textbf{3-state $r_t$ / sticky 0.95} reaches near-$0.5$ ($0.35$--$0.46$) on every ETF. 4-state over-fragments; sign($r_t$) is degenerate on 3 of 7 ETFs.
\item \textbf{$\tau$-shift}: sweep $\tau\in\{0.30,\dots,0.70\}$ in $y_t = \mathbb{1}\{\phi_t^{\mathrm{up}}{>}\tau\}$. Result: moves $\mathrm{label\_mean}$ by $\le 0.03$ on 2-state 2D --- bounded.
\item \textbf{Sticky-floor sweep}: $A_{kk}\,\text{floor}\in\{0.50,\dots,0.99\}$ on 2-state 2D. Result: moves by $\le 0.08$ --- bounded.
\item \textbf{Training-pool shrink}: refit 2-state 2D on last $\{5,10,15,20\}$k bars. Result: no systematic effect; imbalance survives shrinkage.
\end{itemize}

\smallskip
Putting these together, the only lever that actually balances labels in our pipeline is the HMM config choice (3-state $r_t$). But that trades off against signal density (fewer switches).

\textbf{Q2.\enspace Do you target balanced labels, and which lever do you use?} Specifically:
\begin{itemize}
\item Do you accept the balance/signal-density trade-off and pick the HMM config (3-state $r_t$) that gives natural near-$0.5$ labels?
\item Or do you stick with the imbalanced but high-switch 2-state 2D and accept structural imbalance as a feature?
\item \textbf{Should we use $\log(P V)$ as the HMM observation instead of $\log P$?} Our probe shows it flips every ETF from up-majority to up-minority, but the labels then track volume regimes rather than price direction --- which is not what we want to predict. Do you have a use case where labels-on-$\log(P V)$ produces a usable trading signal?
\item Beyond the HMM-config / $\tau$ / sticky / pool / observation levers we measured above, is there an axis-level intervention you would test next?
\end{itemize}

\end{document}
